\section{Wykład 3}
\begin{example}{Mnożenie macierzy}
	\begin{flalign*}
		A,\ B \text{ - macierze } n\times n && & A\cdot B = C = ? && & 
	\end{flalign*}
	\begin{minipage}{0.6\linewidth}
		Najpopularniejszy sposób mnożenia dwóch macierzy $n\times n$ działa w czasie $\BO(n^3)$. 
		Musimy obliczyć $n^2$ wartości i obliczenie każdej zajmuje liniowy czas. 
	\end{minipage}\hfill\begin{minipage}{0.35\linewidth}
		\begin{flalign*}
			\begin{array}{rl}
				B = & \left[\begin{array}{c} \vdots \end{array}\right] \\[2ex]
				A = \left[\begin{array}{cc} \cdots \end{array}\right] 
				& \left[\begin{array}{c} \times \end{array}\right] = C
			\end{array}
		\end{flalign*}
	\end{minipage}
	
	\subsection{Algorytm Strasiera}
	Zakładamy, że $n$ jest potęgą $2$ (dla przejrzystości). Dzielimy $A$ i $B$ na $4$ macierze $\frac{n}{2}\times \frac{n}{2}$:
	
	\begin{flalign*}
		& A = \begin{bNiceArray}{cc}[margin,hvlines]
			A_{11} & A_{12} \\
			A_{21} & A_{22}
		\end{bNiceArray},\  
		B = \begin{bNiceArray}{cc}[margin,hvlines]
		B_{11} & B_{12} \\
		B_{21} & B_{22}
		\end{bNiceArray} && && \\
		%
		& A\cdot B = \begin{bNiceArray}{cc}[margin,hvlines]
			A_{11} & A_{12} \\
			A_{21} & A_{22}
		\end{bNiceArray} \cdot \begin{bNiceArray}{cc}[margin,hvlines]
			B_{11} & B_{12} \\
			B_{21} & B_{22}
		\end{bNiceArray} = \begin{bNiceArray}{cc}[margin,hvlines]
			C_{11} & C_{12} \\
			C_{21} & C_{22}
		\end{bNiceArray} && && \\
		%
		& C_{11} = A_{11}B_{11} + A_{12}B_{21} && && \\
		& C_{12} = A_{11}B_{12} + A_{12}B_{22} && && \\
		& C_{21} = A_{21}B_{11} + A_{21}B_{21} && && \\
		& C_{22} = A_{21}B_{12} + A_{22}B_{22} && && 
	\end{flalign*}
	
	Zamieniliśmy $1$ mnożenie macierzy $n\times n$ na $8$ mnożeń macierzy $\frac{n}{2}\times \frac{n}{2}$ + podział i dodawanie ($\BO(n^2)$). 
	Stąd mamy zależność rekurencyjną: 
	\begin{flalign*}
		T(n) &= 8\cdot T(\frac{n}{2}) + \BT(n^2) && a = 8,\ b = 2,\ k = 2 && \\[1.5ex]
		8 &> 2^2 \Rightarrow T(n) = \BT\left(n^{\log_2 8}\right) = \BT\left(n^3\right) && &&
	\end{flalign*}	

	Strassen zauważył, że jeśli obliczymy
	\begin{flalign*}
		& M_{1} = (A_{12} - A_{22})\cdot (B_{21} + B_{22}) && && \\ 
		& M_{2} = (A_{11} + A_{22})\cdot (B_{11} + B_{22}) && & \text{to } & C_{11} = M_1 + M_2 - M_4 + M_6 \\ 
		& M_{3} = (A_{11} - A_{21})\cdot (B_{11} + B_{12}) && && C_{12} = M_4 + M_5 \\ 
		& M_{4} = (A_{11} + A_{12})\cdot B_{22} && && C_{21} = M_6 + M_7 \\ 
		& M_{5} = A_{11}\cdot (B_{12} - B_{22}) && && C_{22} = M_2 - M_3 + M_5 - M_7 \\ 
		& M_{6} = A_{22}\cdot (B_{21} - B_{11}) && && \\ 
		& M_{7} = (A_{21} + A_{22})\cdot B_{11} && &&
	\end{flalign*}
	
	Zatem
	\begin{flalign*}
		T(n) &= 7\cdot T\left(\frac{n}{2}\right) + \BT(n^2) && a = 7,\ b = 2,\ k = 2 && \\[1.5ex]
		7 &> 2^2 \Rightarrow T(n) = \BT\left(n^{\log_2 7}\right) = \BT\left(n^{2.81}\right) && &&
	\end{flalign*}	
\end{example}

\subsection{Własność optymalnej podstruktury}
Optymalne rozwiązanie problemu jest funkcją optymalnych rozwiązań podproblemów 
(optymalne rozwiązanie można skonstruować z optymalnych rozwiązań podproblemów). 
Żeby metoda \textit{dziel i zwyciężaj} działała dobrze muszą być spełnione dwa warunki (przez problem i algorytm):
\begin{enumerate}
	\item własność optymalnej podstruktury
	\item podproblemy są rozłączne
\end{enumerate}

\subsection{Znajdywanie pary najbliższych punktów}
$Q$ - zbiór $n \geq 2$ punktów na płaszczyźnie
\begin{flalign*}
	p_1 = (x_1, y_1) && p_2 = (x_2, y_2) && d(p_1, p_2) = \sqrt{(x_1 - x_2)^2 + (y_1 - y_2)^2} && && && &&
\end{flalign*}
Problem: znaleźć parę punktów w najmniejszej odległości. Algorytm siłowy (\textit{brute force}), który 
sprawdza wszystkie pary punktów działa w czasie $\BT(n^2)$, bo mamy $\binom{n}{2} = \frac{n(n - 1)}{2} = \BT(n^2) $ par punktów. 
Pokażemy algorytm typu \textit{dziel i zwyciężaj}, który działa w czasie $\BO(n\log n)$. 

W każdym wywołaniu rekurencyjnym wejście składa się z 3 rzeczy: 
zbioru $P\subset Q$ i dwóch tablic $X$ i $Y$ które zawierają wszystkie punkty z $P$. 
W tablicy $X$ punkty są posortowane w porządku niemalejącym względem pierwszej współrzędnej, 
a w tablicy $Y$ po drugiej współrzędnej. 

Opis algorytmu: 
\begin{itemize}
	\item[jeśli $|P| \leq 3$] sprawdzamy wszystkie pary
	\item[jeśli $|P| > 3$] ~ 
	\begin{itemize}
		\item[\underline{Dziel:}] Dzielimy punkty $P$ pionową prostą $l$ na dwa podzbiory $P_L$ i $P_R$ na lewo i na prawo 
		od $l$ tak, że $|P_L| = \left\lceil \frac{P}{2} \right\rceil,\ |P_R| = \left\lfloor \frac{P}{2} \right\rfloor$. 
		Dzielimy tablicę $X$ na $X_L$ i $X_R$ zawierające punkty odpowiednio $P_L$ i $P_R$ posortowane w porządku 
		niemalejącym względem pierwszej współrzędnej. Podobnie dzielimy $Y$ na $Y_L$ i $Y_R$ zawierające 
		odpowiednio punkty z $P_L$ i $P_R$ posortowane w porządku niemalejącym po drugiej współrzędnej. 
		
		\item[\underline{Zwyciężaj:}] Wywołujemy rekurencyjnie algorytm dla wejścia $P_L,\ X_L,\ Y_L$ i $P_R,\ X_R,\ Y_R$. 
		Algorytm znajduje pary najbliższych punktów w $P_L$ i $P_R$. Niech $\delta_L$ i $\delta_R$ będą 
		najmniejszymi odległościami znalezionymi przez algorytm dla odpowiednio $P_L$ i $P_R$. 
		Niech $\delta$ będzie mniejszą z nich: $\delta = \min(\delta_L,\delta_R)$.
		
		\item[\underline{Połącz:}] Para najbliższych punktów z $P$ składa się z dwóch punktów z $P_L$ lub z dwóch punktów 
		z $P_R$ lub jednego punkty z $P_L$ i jednego z $P_R$. 
		
		Musimy sprawdzić czy istnieje para 3-go typu w odległości mniejszej niż $\delta$. 
		Jeśli taka para istnieje, wtedy jej punkty są w pionowym pasie o szerokości $2\delta$ 
		wokół $l$. 
		
		W celu sprawdzenia czy taka para istnieje wykonujemy następujące czynności:
		\begin{enumerate}
			\item Tworzymy tablicę $Y'$ otrzymaną z $Y$ poprzez usunięcie punktów spoza pasa. $Y'$ jest 
			posortowany w porządku niemalejącym względem drugiej współrzędnej. 
			\item Dla każedego punktu $p$ z $Y'$ sprawdzamy odległość między $p$, a $p'$ dla $7$ kolejnych 
			punktów $p'$ z $Y'$. Algorytm oblicza najmniejszą odległość $\delta'$ i zapamiętuje ją i odpowiednią 
			parę punktów. 
			\item Jeśli $\delta' < \delta$ to pionowy pas zawiera parę punktów w mniejszej odległości niż pary znalezione 
			w wywołaniach rekurencyjnych. Algorytm zwraca $\delta'$ i odpowiednią parę punktów. Wpp. zwraca $\delta$ 
			i odpowiednią parę punktów. 
		\end{enumerate}
	\end{itemize}
\end{itemize}

\begin{proof}[Dowód poprawności]
	Jedyną rzeczą, która nie jest oczywista jest fakt, że sprawdzamy tylko $7$ następnych punktów po $p$ w $Y'$ w fazie połącz. 
	Przypuśćmy, że $p, q$ są parą punktów w najmniejszej odległości w $P$ i $\delta(p,q) = \delta'$. Niech 
	$p = (x_p, y_p),\ q = (x_q, y_q)$. Bez straty ogólności załóżmy $y_p \leq y_q$. Pokażemy, że algorytm znajdzie 
	tę parę punktów. Przypuśćmy, że nie. Wtedy istnieje co najmniej $7$ punktów pomiędzy $p$ i $q$ w $Y'$. 
	Niech $r = (x_r, y_r)$ będzie jednym z nich. Skoro $d(p,q) < \delta$ to 
	$y_q - y_p = |y_q - y_p| \leq \sqrt{(x_p - x_q)^2 + (y_p - y_q)^2} = \delta(p, q) \leq \delta$. 
	Ponadto $y_r$ musi być między $y_p$ i $y_q$: $y_p \leq y_r \leq y_q$, więc $y_r - y_p \leq y_q - y_p < \delta$ 
	zatem co najmniej 9 punktów z $P$ w prostokącie $2\delta\times\delta$ ograniczonym poziomymi liniami 
	$y = y_p,\ y = y_p + \delta$. 
	
	Wtedy w lewym lub prawym kwadracie $\delta\times\delta$ jest przynajmniej $5$ punktów z $P$. Przypuśćmy, 
	że to lewy kwadrat. Zatem istnieje kwadrat $\delta\times\delta$ z $5$ punktami z $P$ w $P_L$. 
	Ale wtedy istnieją dwa punkty w $P_L$ w odległości mniejszej niż $\delta$, bo nie da się umieścić 
	w kwadracie $\delta\times\delta$ $5$ punktów dla których odległość między dowolnymi dwoma 
	jest mniejsza niż $\delta$. Sprzeczność (bo najmniejsza odległość w $P_L$ to $\delta_L \geq \delta$). 
	
	Zatem wystarczy sprawdzić tylko $7$ następnych punktów po $p$ w $Y'$. 
\end{proof}

\paragraph{Analiza złożoności czasowej}
Na początku musimy posortować tablicę $X$ i $Y$. To zajmuje $\BO(n\log n)$. Jeśli $X$ i $Y$ są już posortowane 
to podział $X$ na posortowane $X_L$ i $X_L$ oraz $Y$ analogicznie zajmie $\BO(n)$. 

\begin{alg}{???}{alg:closestpoints:one}
	\State \Length{$Y_L$} \Assign \Length{$Y_R$} \Assign $0$
	\For{$i$ \Assign $1$ \To \Length{$Y$}}
		\If{$Y(i)\in P_L$}
			\State \Length{$Y_L$} \Assign \Length{$Y_L$} + 1
			\State $Y_L($\Length{$Y_L$}$)$ \Assign $Y(i)$
		\Else
			\State \Length{$Y_R$} \Assign \Length{$Y_R$} + 1
			\State $Y_R($\Length{$Y_R$}$)$ \Assign $Y(i)$
		\EndIf
	\EndFor
\end{alg}

Mamy dwa wywołania rekurencyjne dla $n/2$ punktów i fazę połącz. Tablicę $Y'$ można otrzymać z $Y$ w czasie $\BO(n)$ 
tak, że $Y'$ jest posortowana. Niech $x = x_L$ będzie równanie prostej $l$. 

\begin{alg}{???}{alg:closestpoints:two}
	\State \Length{$Y'$} \Assign $0$
	\For{$i$\Assign $1$ \To \Length{$Y$}}
		\If{$x_L - \delta \leq Y(i).x \leq x_L + \delta$}
			\State \Length{$Y'$} \Assign \Length{$Y'$} + 1
			\State $Y'($\Length{$Y'$}$)$ \Assign $Y(i)$
		\EndIf
	\EndFor
\end{alg}

Skoro dla każdego z $\BO(n)$ punktów z $Y'$ sprawdzamy odległość tylko co najwyżej $7$ punktów, to najmniejsza odległość 
$\delta'$ punktów z $Y'$ znajdziemy w czasie $\BO(n)$. 

Oznaczmy złożoność algorytmu bez wstępnego sortowania przez $T(n)$. Mamy:
\begin{flalign*}
	T(n) &= 2\cdot T\left(\frac{n}{2}\right) + \BT(n^1) && a = 2,\ b = 2,\ k = 1 && \\[1.5ex]
	2 &= 2^1 \Rightarrow T(n) = \BT\left(n\log n\right) && &&
\end{flalign*}	

 
  
